\chapter{Sviluppo dell'estensione}
\label{chap:sviluppo-estensione}

Questo capitolo racconta lo stage dal punto di vista del metodo di lavoro adottato, del contesto tecnico e normativo necessario a inquadrare il problema, dell'analisi dei requisiti, delle scelte di progettazione e implementazione, e infine della verifica e dei risultati raggiunti.

\section{Metodo di lavoro}
\label{sec:metodo-lavoro}

% TODO: pianificazione delle attività, modalità di interazione con il tutor aziendale,
% cadenza delle revisioni di progresso, strumenti usati per tracciare requisiti e progressi.


\section{Il quadro normativo e tecnico di riferimento}
\label{sec:quadro-normativo}

\subsection{Cos'è l'accessibilità e qual è la sua finalità}
\label{subsec:cose-accessibilita}

L'accessibilità digitale può essere definita, in primo luogo, richiamando la definizione fornita dalla Legge del 9 gennaio 2004, n.~4, comunemente nota come \enquote{legge Stanca}, la prima legge italiana dedicata al tema: \enquote{la capacità dei sistemi informatici ivi inclusi i siti web e le applicazioni mobili, nelle forme e nei limiti consentiti dalle conoscenze tecnologiche, di erogare servizi e fornire informazioni fruibili, senza discriminazioni, anche da parte di coloro che a causa di disabilità necessitano di tecnologie assistive o configurazioni particolari}\footcite[art.~2, comma~1, lett.~a]{legge:stanca}. La stessa legge riconduce esplicitamente tale obiettivo al principio di uguaglianza sostanziale sancito dall'art.~3 della Costituzione, stabilendo che l'accesso delle persone con disabilità ai servizi informatici sia tutelato proprio \enquote{in ottemperanza} a tale principio\footcite[art.~1, comma~2]{legge:stanca}. Una definizione più generale, di matrice tecnica anziché giuridica, è quella adottata dalla norma europea EN~301~549, che riprende a sua volta la norma ISO~9241-11: l'accessibilità è la \enquote{misura in cui i prodotti, i sistemi, i servizi, gli ambienti e le strutture possono essere utilizzati da persone di una popolazione con la più ampia gamma di necessità degli utenti, caratteristiche e capacità, per raggiungere un obiettivo identificato in un contesto identificato di utilizzo}\footcite[cap.~3, voce \enquote{accessibilità}, da ISO~9241-11:2018]{norma:en301549}. Le due definizioni si completano: la prima pone in evidenza la necessità di garantire l'accesso alle persone con disabilità, anche attraverso tecnologie assistive o configurazioni particolari, e colloca tale obiettivo nel quadro del principio di uguaglianza; la seconda descrive l'accessibilità in termini più ampi, come proprietà di un sistema rispetto a una popolazione di utenti eterogenea per caratteristiche e capacità, di cui le persone con disabilità sono una componente e non la totalità.

Questa prospettiva più ampia emerge con particolare chiarezza nella struttura della norma EN~301~549, in particolare nel capitolo~4, dedicato alle prestazioni funzionali che un sistema accessibile deve garantire\footcite[cap.~4 \enquote{Prestazioni funzionali}]{norma:en301549}. Per ciascuno dei dieci criteri individuati, la norma specifica la categoria di utenti per cui il requisito risulta essenziale e, al contempo, evidenzia come la medesima caratteristica possa apportare benefici anche in altre situazioni. Ad esempio, per l'utilizzo senza vista, la norma richiede che l'ICT fornisca una modalità operativa che non richieda la vista e precisa: \enquote{Questo è essenziale per gli utenti senza vista e avvantaggia molti altri utenti in diverse situazioni}\footcite[§~4.2.1 \enquote{Utilizzo senza vista}]{norma:en301549}. Una formulazione pressoché identica ricorre nelle successive sezioni, ciascuna dedicata a una diversa differenza di abilità, adattando di volta in volta il riferimento alla specifica condizione considerata.

Il principio è facilmente osservabile anche in caratteristiche di uso comune: i sottotitoli sono essenziali per chi presenta una riduzione parziale o una perdita totale della capacità di percepire i suoni, ma risultano utili anche a chi guarda un video in un ambiente rumoroso o senza audio disponibile; un'alternativa testuale a un'immagine è essenziale per chi utilizza uno screen reader, ma può risultare utile anche quando l'immagine non viene caricata correttamente o quando la connessione di rete è lenta; un adeguato contrasto tra testo e sfondo, infine, è particolarmente importante per alcune differenze di abilità visive, ma può migliorare la leggibilità anche in condizioni ambientali sfavorevoli, ad esempio quando lo schermo viene utilizzato sotto una forte illuminazione. L'accessibilità non consiste, quindi, soltanto nel rimuovere ostacoli per una specifica categoria di utenti, ma nel progettare sistemi capaci di adattarsi a una varietà di caratteristiche e condizioni d'uso.

Questa impostazione si riflette anche nel modo in cui viene delineata la platea dei beneficiari dei requisiti di accessibilità. La direttiva europea sui requisiti di accessibilità dei prodotti e dei servizi, il cosiddetto European Accessibility Act (EAA), richiama la definizione di disabilità della Convenzione ONU sui diritti delle persone con disabilità, riferita alle persone che presentano minorazioni fisiche, mentali, intellettive o sensoriali che, in interazione con barriere di varia natura, possono ostacolarne la piena partecipazione alla società\footcite[considerando~3, che richiama la Convenzione delle Nazioni Unite sui diritti delle persone con disabilità (UNCRPD), ratificata dall'Unione europea il 21 gennaio 2011]{direttiva:eaa}. La stessa direttiva riconosce, tuttavia, che i requisiti di accessibilità possono apportare benefici anche a \enquote{altre persone con limitazioni funzionali, tra cui, ad esempio, le persone anziane, le donne in gravidanza o chi si trova a viaggiare con un bagaglio ingombrante}\footcite[considerando~4]{direttiva:eaa}. 

La necessità di considerare condizioni d'uso differenti porta quindi a superare una concezione dell'accessibilità legata esclusivamente a condizioni permanenti e diagnosticate. Una barriera può infatti essere associata a una condizione permanente, come una lieve ipovisione o una difficoltà di lettura; a una condizione temporanea, come un braccio ingessato che impedisce l'utilizzo del mouse; oppure a una condizione situazionale, come la difficoltà di leggere uno schermo esposto alla luce diretta del sole, di ascoltare un contenuto audio in un ambiente rumoroso o di utilizzare le mani mentre si tiene in braccio un bambino. In tutti questi casi, la barriera non risiede necessariamente nella persona, ma può emergere dall'interazione tra le caratteristiche dell'utente, il contesto e le modalità con cui il sistema è stato progettato.

La distinzione tra condizioni permanenti, temporanee e situazionali costituisce pertanto uno dei criteri concettuali che guidano le scelte progettuali del presente lavoro. Un'estensione finalizzata alla verifica dell'accessibilità di una pagina web risulta tanto più significativa quanto più i criteri selezionati sono in grado di intercettare barriere che possono manifestarsi trasversalmente nelle diverse condizioni d'uso, anziché rivolgersi esclusivamente a una singola categoria di utenti. L'obiettivo dell'accessibilità diventa così non soltanto garantire l'utilizzabilità del sistema da parte di chi incontra una specifica barriera, ma contribuire alla realizzazione di interfacce più robuste e fruibili in una maggiore varietà di situazioni.

\subsection{Cosa deve garantire il web: i quattro principi dell'accessibilità}
\label{subsec:quattro-principi}

Se l'accessibilità riguarda una platea di utenti così eterogenea, è necessario un criterio unificante che permetta di valutare in modo oggettivo se un contenuto digitale sia o meno fruibile, indipendentemente dalla natura specifica della barriera che un utente incontra. Questo criterio è stato elaborato dal World Wide Web Consortium (W3C) attraverso le Web Content Accessibility Guidelines (\gls{wcag}), lo standard tecnico internazionale che struttura l'accessibilità attorno a quattro principi, noti con l'acronimo POUR: Percepibile, Operabile (o \enquote{Utilizzabile}, nella traduzione italiana), Comprensibile, Robusto.

Il documento ufficiale che ne illustra il significato definisce ciascun principio in termini concreti e verificabili nella pratica\footcite[§ \enquote{Understanding the Four Principles of Accessibility}]{w3c:understanding-intro}. Percepibile perché \enquote{le informazioni e i componenti dell'interfaccia utente devono essere presentabili agli utenti in modi che essi possano percepire}: le informazioni devono raggiungere l'utente attraverso almeno un canale sensoriale compatibile con le sue capacità, perché un contenuto solo visivo esclude chi non vede e un contenuto solo audio esclude chi non sente. Operabile perché \enquote{i componenti dell'interfaccia utente e la navigazione devono essere operabili}: l'interazione non può dipendere da un'unica modalità di input, perché un sito navigabile solo col mouse esclude chi usa la tastiera o un dispositivo di puntamento alternativo. Comprensibile perché \enquote{le informazioni e il funzionamento dell'interfaccia utente devono essere comprensibili}: il contenuto e il suo funzionamento non possono superare la capacità di comprensione dell'utente, requisito che intercetta tanto le disabilità cognitive quanto, più banalmente, l'utente distratto o in contesto di fretta. Robusto perché \enquote{il contenuto deve essere sufficientemente solido da poter essere interpretato in modo affidabile da un'ampia varietà di programmi utente, incluse le tecnologie assistive}: il contenuto deve restare interpretabile correttamente anche da tecnologie diverse, comprese le tecnologie assistive come gli screen reader, requisito che sposta parte della responsabilità dal solo utente alla qualità tecnica del codice, e che deve restare valido anche al progredire delle tecnologie stesse\footcite[\enquote{If any of these are not true, users with disabilities will not be able to use the web.}]{w3c:understanding-intro}. Se anche uno solo di questi quattro principi non è rispettato, conclude il documento, gli utenti con disabilità non sono in grado di usare il web.

Dai quattro principi discendono, secondo la struttura ufficiale delle WCAG, tredici linee guida (\emph{guidelines}), obiettivi generali non direttamente verificabili, e, per ciascuna di esse, uno o più criteri di successo (\emph{success criteria}): affermazioni verificabili, formulate in modo da poter essere valutate come vere o false quando applicate a un contenuto specifico, indipendenti dalla tecnologia utilizzata\footcite[§~0.2 \enquote{WCAG 2 Livelli di orientamento}]{w3c:wcag21}. È a questo terzo livello, i criteri di successo, che si aggiunge la classificazione per livello di conformità (A, AA, AAA) discussa nella sezione~\ref{subsec:livello-aa}. Un quarto livello, informativo e non normativo, raccoglie infine le tecniche sufficienti e consigliate per soddisfare ciascun criterio, documentate dal W3C in un catalogo consultabile e filtrabile\footcite{w3c:quickref}.

La cosa rilevante, ai fini di questa sezione, è che questi quattro principi non restano un'elaborazione tecnica confinata agli ambienti W3C: si ritrovano, pressoché identici nella formulazione, in ogni livello della normativa ricostruita nella sezione successiva. La legge Stanca, nella versione aggiornata dopo il recepimento della direttiva europea sui siti web delle pubbliche amministrazioni, dichiara accessibili i siti e le applicazioni mobili che siano \enquote{percepibili, utilizzabili, comprensibili e solidi}\footcite[art.~3-bis, comma~1]{legge:stanca}. La norma tecnica europea EN~301~549, che è il riferimento tecnico condiviso da tutta la normativa successiva, dedica un intero capitolo alle \enquote{prestazioni funzionali} che declinano gli stessi principi in dieci criteri operativi (utilizzo senza vista, con vista limitata, senza percezione del colore, senza udito, e così via)\footcite[cap.~4 \enquote{Prestazioni funzionali}, punti 4.2.1-4.2.13]{norma:en301549}. Il decreto legislativo che recepisce l'EAA in Italia riproduce la medesima struttura nell'Allegato~I, Sezione~VII. Le stesse linee guida AGID più recenti, nel presentare i \enquote{quattro pilastri dell'accessibilità}, li riportano nella formulazione ormai canonica\footcite[§~6.1]{agid:lg-82-2022}. Il fatto che principi identici attraversino trent'anni di elaborazione tecnica e livelli normativi diversi (internazionale, europeo, nazionale) segnala che non si tratta di una convenzione arbitraria, ma di una tassonomia che intercetta in modo esaustivo le modalità con cui una barriera digitale può manifestarsi. È per questo che, come si vedrà nella sezione~\ref{sec:selezione-criteri}, la selezione dei criteri \gls{wcag} oggetto di verifica automatica nel presente lavoro è stata condotta avendo cura di rappresentare, per quanto possibile, tutti e quattro i principi, e non solo quello --- la percepibilità --- più immediatamente associato all'idea comune di \enquote{disabilità visiva}.


\subsection{Come questi principi diventano obbligo giuridico: la stratificazione delle fonti}
\label{subsec:stratificazione-fonti}

Un conto è disporre di un principio tecnico condiviso, un altro è tradurlo in un obbligo giuridico verificabile ed esigibile. Il percorso che ha condotto le \gls{wcag} da raccomandazione volontaria del W3C a standard richiamato da norme cogenti è avvenuto per stratificazione progressiva, su tre livelli che è utile ricostruire in sequenza, perché ciascuno risponde a un'esigenza diversa e ciascuno resta tuttora in vigore, sovrapposto agli altri.

\subsubsection{Il livello italiano pionieristico: la legge Stanca (2004)}

L'Italia è stata tra i primi paesi al mondo a dotarsi di una legge organica sull'accessibilità informatica, anticipando di oltre un decennio l'intervento dell'Unione europea. La legge 9 gennaio 2004, n.~4 nasce però con un ambito soggettivo circoscritto: si rivolge alle pubbliche amministrazioni, agli enti pubblici economici, alle aziende concessionarie di servizi pubblici e --- solo in una fase successiva --- ai soggetti privati di grandi dimensioni (con una soglia di fatturato medio superiore a cinquecento milioni di euro)\footcite[art.~3, commi~1 e 1-bis]{legge:stanca}. È una legge \enquote{a delega tecnica aperta}: l'art.~11 non definisce in proprio i requisiti tecnici di accessibilità, ma affida all'Agenzia per l'Italia Digitale (\gls{agid}) il compito di emanarli tramite linee guida, da aggiornare periodicamente e da allineare agli standard tecnici europei e internazionali. È una scelta di architettura normativa che si rivelerà decisiva: separando il principio giuridico (l'obbligo di accessibilità) dal dettaglio tecnico (i singoli requisiti da soddisfare), la legge può restare stabile nel tempo mentre lo standard tecnico sottostante evolve --- dalle \gls{wcag} 1.0 alle attuali 2.2 --- senza richiedere ogni volta un intervento del legislatore.

L'estensione dell'obbligo ai soggetti privati di grandi dimensioni (art.~3, comma~1-bis) ha portato \gls{agid} a emanare, accanto alle linee guida generali, un documento distinto dedicato specificamente a questi soggetti\footcite{agid:lg-stanca-1bis}. Il confronto tra i due testi è istruttivo, perché mostra che l'estensione della platea degli obbligati non ha comportato un'estensione proporzionale dell'apparato di controllo. Sul piano dei requisiti tecnici, i due regimi coincidono: stessa norma EN~301~549, stessi capitoli di riferimento per hardware, web, documenti non web, software, applicazioni mobili e documentazione, stessa soglia di conformità \gls{wcag} 2.1 livello AA. A cambiare in modo sostanziale è tutto ciò che riguarda la verifica e il controllo di quell'obbligo. Per le pubbliche amministrazioni, le linee guida disciplinano nel dettaglio una procedura di verifica soggettiva articolata in quattro fasi --- dall'analisi di un esperto di fattori umani fino alla costituzione di un gruppo di valutazione che coinvolge persone con diverse disabilità --- e la rendono obbligatoria per le forniture sopra soglia comunitaria; prevedono inoltre un intero capitolo dedicato alla metodologia di monitoraggio, con cui \gls{agid} verifica annualmente, su un campione di siti calcolato in proporzione alla popolazione, la conformità effettiva degli enti pubblici, riferendone alla Commissione europea con cadenza triennale\footcite[cap.~3 e cap.~5]{agid:lg-stanca}. Per i soggetti privati di cui al comma~1-bis, invece, la verifica soggettiva resta una possibilità suggerita --- \enquote{è possibile effettuare test di usabilità [...] con il coinvolgimento anche di persone con disabilità} --- senza una metodologia prescritta né un vincolo di obbligatorietà legato a soglie, e non è previsto alcun capitolo equivalente sul monitoraggio periodico da parte di \gls{agid}\footcite[§~3.2.2]{agid:lg-stanca-1bis}. Il controllo, per questi soggetti, resta di natura reattiva: si attiva a seguito di una segnalazione dell'utente, e le segnalazioni ricevute concorrono, semmai, a orientare la pianificazione di verifiche d'ufficio successive, piuttosto che alimentare un monitoraggio sistematico e programmato come quello riservato alla pubblica amministrazione.

Questa differenza ha una conseguenza pratica rilevante: per i siti della pubblica amministrazione, la conformità reale ai requisiti di accessibilità viene verificata periodicamente da un soggetto terzo (\gls{agid}), con una metodologia definita e un campione statisticamente rappresentativo; per i grandi soggetti privati obbligati (art.~3, comma~1-bis), questo controllo esterno sistematico non è previsto, e la verifica dipende dall'iniziativa spontanea dell'ente o da un'eventuale segnalazione.

\subsubsection{Il livello europeo: dalla Direttiva 2016/2102 all'European Accessibility Act}

L'Unione europea interviene in due tempi. Con la direttiva 2016/2102 impone agli enti pubblici degli Stati membri di rendere accessibili i propri siti web e applicazioni mobili, individuando nella norma tecnica armonizzata EN~301~549 (elaborata dagli organismi di normazione europei CEN, CENELEC ed ETSI) lo strumento con cui verificare la conformità: è questa norma tecnica, e non la direttiva stessa, a stabilire che la conformità al livello AA delle \gls{wcag} equivale al soddisfacimento dei requisiti giuridici europei\footcite[§~9.0: \enquote{La conformità con il livello AA delle W3C Web Content Accessibility Guidelines (WCAG 2.1) è equivalente alla conformità con tutti i punti da 9.1 a 9.4 e ai requisiti di conformità di cui al punto 9.6 del presente documento}]{norma:en301549}. La norma EN~301~549 funge quindi da cerniera tecnica tra due mondi che altrimenti non comunicherebbero direttamente: da un lato il diritto, che parla il linguaggio degli obblighi e delle sanzioni; dall'altro lo standard W3C, che parla il linguaggio dei criteri di successo verificabili.

Il secondo intervento, più ampio, è la direttiva (UE) 2019/882, il già citato European Accessibility Act, adottata nel 2019 ma con efficacia degli obblighi rinviata al 28 giugno 2025 per dare agli Stati membri e agli operatori economici il tempo di adeguarsi. Se la direttiva del 2016 riguardava solo il settore pubblico, l'EAA estende gli obblighi di accessibilità a una lista mirata di prodotti e servizi privati ad alto impatto sulla vita quotidiana: sistemi operativi e hardware di consumo, sportelli self-service (come sportelli bancomat e biglietterie automatiche), servizi di comunicazione elettronica, servizi bancari per consumatori, commercio elettronico, e-book e relativi lettori, trasporto passeggeri\footcite[art.~2]{direttiva:eaa}\footcite[art.~1, commi~2-3]{dlgs:82-2022}. È la presa d'atto che l'accessibilità digitale non è più (se mai lo è stata) un problema circoscritto ai portali della pubblica amministrazione: riguarda ogni servizio digitale con cui un cittadino interagisce quotidianamente, dalla app della propria banca al sito di un e-commerce.

\subsubsection{Il recepimento italiano: il D.Lgs.~82/2022 e l'aggiornamento della legge Stanca}

Sotto il profilo dei soggetti obbligati, il D.Lgs.~82/2022 segue una logica diversa da quella della legge Stanca. Mentre quest'ultima individua i soggetti privati obbligati attraverso una soglia dimensionale --- il fatturato medio superiore a cinquecento milioni di euro, che esclude in blocco tutte le imprese di dimensioni minori --- il decreto del 2022 delimita il proprio ambito di applicazione non in base alla dimensione dell'impresa, ma in base al tipo di prodotto o servizio offerto, elencato tassativamente dall'art.~1, commi~2 e~3\footcite[art.~1, commi~2-3]{dlgs:82-2022}. Ne consegue che l'obbligo si estende, in linea di principio, anche alle microimprese, purché immettano sul mercato uno dei prodotti elencati (ad esempio hardware di consumo o terminali self-service): il decreto esenta espressamente le sole microimprese che forniscono uno dei servizi elencati dal comma~3, non quelle che producono uno dei prodotti di cui al comma~2\footcite[art.~3, comma~3: \enquote{Le microimprese che forniscono servizi sono esentate dall'osservanza dei requisiti di accessibilità di cui al comma 2}]{dlgs:82-2022}. Il perimetro dei soggetti raggiunti dal D.Lgs.~82/2022 è quindi, per questa componente, potenzialmente più fine e più esteso verso il basso di quello della legge Stanca, benché limitato alle sole categorie di prodotto e servizio espressamente elencate dalla norma.

L'Italia recepisce l'EAA con il decreto legislativo 27 maggio 2022, n.~82\footnote{Da non confondere con il D.Lgs.~7 marzo 2005, n.~82 (Codice dell'Amministrazione Digitale)~\autocite{dlgs:82-2005}, che ha numero identico ma anno di emanazione diverso: nel prosieguo del lavoro si farà sempre riferimento esplicito all'anno per evitare ambiguità.}, le cui disposizioni sono divenute efficaci dal 28 giugno 2025. Il decreto adotta la medesima architettura già collaudata dalla legge Stanca: definisce i requisiti di accessibilità per macro-categoria di prodotto o servizio (Allegato~I), ma per la loro traduzione operativa rinvia nuovamente alla norma tecnica EN~301~549 e a linee guida attuative emanate da \gls{agid}, pubblicate in attuazione dell'art.~21 del decreto stesso. In parallelo, la legge Stanca originaria non è stata abrogata ma progressivamente aggiornata (da ultimo con il D.Lgs.~10 agosto 2018, n.~106, che ne ha innovato gli articoli sui principi generali, sulla dichiarazione di accessibilità e sul monitoraggio), così che oggi coesistono due corpi normativi paralleli --- uno storicamente centrato sulla pubblica amministrazione (con un'estensione ai grandi soggetti privati basata sul fatturato), l'altro su un elenco di prodotti e servizi privati definiti per tipologia, indipendentemente dalla dimensione d'impresa salvo le eccezioni per microimprese di servizi --- ma entrambi ancorati allo stesso standard tecnico di riferimento.

È proprio questa convergenza tecnica, al di là della complessità delle fonti giuridiche e dei diversi criteri con cui ciascuna di esse individua i propri soggetti obbligati, il dato più rilevante ai fini del presente lavoro: qualunque sia il soggetto obbligato e qualunque sia la fonte formale dell'obbligo, il contenuto tecnico del requisito di accessibilità per un sito web resta lo stesso --- conformità \gls{wcag} 2.1 livello AA, come stabilito dalla norma EN~301~549 e confermato esplicitamente dalle linee guida \gls{agid} sia per la legge Stanca sia per il D.Lgs.~82/2022\footcite[§~2.2]{agid:lg-stanca}\footcite[§~8.5]{agid:lg-82-2022}.

\subsubsection{Un catalogo tecnico stabile, un'applicazione normativa variabile}

I paragrafi precedenti permettono di distinguere due piani che è utile tenere separati. Da un lato c'è il catalogo tecnico dei criteri \gls{wcag}: un elenco fisso e tecnologicamente neutro, definito una sola volta dal W3C e richiamato senza modifiche da ogni fonte normativa esaminata --- lo stesso criterio 1.4.3 \enquote{Contrasto (minimo)}, ad esempio, ha lo stesso contenuto tecnico per un sito della pubblica amministrazione, per l'app di una banca o per un e-commerce qualsiasi. Dall'altro lato c'è l'applicazione concreta di quel catalogo a un caso specifico, che invece non è affatto stabile: dipende dall'intreccio di più variabili che ciascuna fonte normativa definisce autonomamente, e che possono sovrapporsi solo in parte.

Una prima variabile è chi gestisce il sito. Un ente pubblico è sempre soggetto alla legge Stanca; un soggetto privato lo è solo se supera la soglia di fatturato di cinquecento milioni di euro (art.~3, comma~1-bis) oppure se rientra in una delle categorie di prodotto o servizio elencate dal D.Lgs.~82/2022, indipendentemente dal fatturato, salvo l'eccezione delle microimprese di servizi discussa poco sopra. Uno stesso sito privato potrebbe quindi essere obbligato da una sola delle due fonti, da entrambe, oppure da nessuna, a seconda di come si incrociano dimensione dell'impresa e tipologia di servizio offerto.

Una seconda variabile è il tipo di contenuto interessato. Come si è visto discutendo la norma EN~301~549 e l'Allegato~I del D.Lgs.~82/2022, i criteri tecnici applicabili non sono gli stessi per una pagina web, per un documento scaricabile e per un software: cambiano il capitolo di riferimento della norma tecnica (9, 10 o 11) e la sezione dell'Allegato~I richiamata (III o IV, a seconda che si tratti di un requisito generale di servizio o di un requisito specifico per una particolare tipologia di servizio, come i servizi bancari o l'e-commerce). Un sito può quindi rispettare pienamente i requisiti previsti per le proprie pagine web e, allo stesso tempo, non essere conforme per un documento PDF allegato, perché i due contenuti sono valutati secondo capitoli diversi della stessa norma.

Una terza variabile è se sia richiesta una conformità piena o solo parziale. Entrambe le fonti prevedono meccanismi di deroga: la legge Stanca ammette l'esclusione di singoli requisiti quando comportino un \enquote{onere sproporzionato}, nei casi tassativamente elencati dall'art.~3-ter e ulteriormente specificati dalle linee guida \gls{agid} (che, ad esempio, escludono a priori contenuti come le riproduzioni di beni del patrimonio storico-culturale o le dirette in streaming)\footcite[§~6.1.4]{agid:lg-stanca}; il D.Lgs.~82/2022 prevede una clausola equivalente, articolata secondo criteri di proporzionalità tra costi dell'adeguamento e beneficio per gli utenti. In entrambi i casi, quindi, il fatto che un criterio compaia nel catalogo \gls{wcag} non implica automaticamente che il suo soddisfacimento sia dovuto senza eccezioni: l'obbligo reale è sempre il risultato di un catalogo di per sé fisso, filtrato attraverso le condizioni di deroga previste dalla fonte normativa applicabile al caso concreto.

A queste variabili si aggiunge, come illustrato poco sopra, una quarta dimensione che riguarda non il contenuto dell'obbligo ma la sua verificabilità: chi controlla che l'obbligo sia effettivamente rispettato, e con quale sistematicità, cambia sensibilmente a seconda che si tratti di un ente pubblico, di un grande soggetto privato o di un sito che non rientra in alcuna delle due soglie. È dalla somma di queste quattro variabili --- soggetto obbligato, tipo di contenuto, ampiezza della conformità richiesta, e sistematicità del controllo --- che deriva la complessità applicativa dell'accessibilità digitale in Italia: i criteri tecnici, in sé, sono chiari e stabili; la loro effettiva esigibilità per un sito determinato, molto meno.


\subsection{Perché il livello AA e perché solo alcuni dei criteri WCAG}
\label{subsec:livello-aa}

Le \gls{wcag} articolano ogni criterio di successo su tre livelli di conformità crescente --- A, AA, AAA --- che non rappresentano gradi di importanza ma un compromesso via via più stringente tra copertura delle disabilità e fattibilità tecnica/economica dell'implementazione. Il livello AAA, pur desiderabile, non è richiesto come requisito normativo generale: lo stesso W3C sconsiglia di imporlo come obbligo per un intero sito, perché alcuni suoi criteri non sono soddisfacibili da ogni tipologia di contenuto\footcite[§~9.0, nota~4]{norma:en301549}. È per questo che tutta la normativa ricostruita nella sezione precedente --- europea e italiana --- converge sul livello AA come soglia di conformità legale: è il livello richiesto dalla norma EN~301~549 (punti 9.1-9.4 più i requisiti di conformità del punto~9.6), ed è quindi, di fatto, l'unico livello dotato di valore normativo cogente nel contesto che qui interessa. Va precisato che il riferimento è, allo stato attuale, alla versione 2.1 delle \gls{wcag} (o, per i siti sviluppati prima dell'entrata in vigore delle linee guida \gls{agid}, alla precedente versione 2.0): la norma EN~301~549 è infatti attualmente in fase di aggiornamento per allinearsi anche alla versione 2.2 delle \gls{wcag}, pubblicata dal W3C nell'ottobre 2023 e già disponibile in traduzione italiana ufficiale\footcite[§~8.5 e §~2 (Riferimenti norme tecniche)]{agid:lg-82-2022}, per cui nel prosieguo del lavoro si farà riferimento alla versione 2.1, tuttora quella giuridicamente vincolante.

All'interno del livello AA, i criteri di successo applicabili --- considerando anche i vuoti numerici lasciati per allineamento con i livelli superiori --- sono raggruppati dalla norma EN~301~549 in modo identico per tre contesti distinti: le pagine web (capitolo~9), i documenti non web come i PDF scaricabili (capitolo~10) e, per riflesso, la documentazione e i servizi di supporto (capitolo~12), che rinviano agli stessi requisiti dei primi due. Contando i criteri di successo elencati nel testo ufficiale delle WCAG 2.1, quelli di livello A sono 30 e quelli di livello AA sono 20, per un totale di 50 criteri obbligatori a livello AA (sugli 78 complessivi, includendo i 28 di livello AAA)\footcite{w3c:quickref}. È questo l'insieme di riferimento --- i 50 criteri \gls{wcag} 2.1 AA applicabili al web, ai documenti scaricabili e alla documentazione di supporto --- da cui il presente lavoro seleziona il proprio sottoinsieme operativo. Va segnalato, per completezza, che alcune fonti secondarie e materiali di lavoro preliminari a questo studio riportavano talvolta la cifra di 51 criteri: un controllo incrociato con il testo normativo ufficiale delle WCAG 2.1 non ha permesso di individuare alcuna fonte primaria a sostegno di tale numero, né esso risulta corroborato dalla norma EN~301~549, che si limita a rinviare \enquote{a tutti i punti da 9.1 a 9.4} delle WCAG senza indicarne il totale numerico. È possibile che la cifra derivi da una confusione con il totale dei criteri di livello A/AA delle WCAG 2.2 (55, per effetto dei sei nuovi criteri introdotti e del ritiro del criterio 4.1.1), oppure da un conteggio manuale impreciso. Nel prosieguo del lavoro si farà quindi riferimento al numero di 50, verificato direttamente sul testo ufficiale delle WCAG 2.1.

La scelta di non verificare la totalità dei criteri, ma un numero ristretto selezionato, nasce da un vincolo pratico e da un'opportunità tecnica che caratterizzano questo lavoro. Il vincolo pratico è che non tutti i criteri \gls{wcag} sono ugualmente verificabili in modo affidabile da un sistema automatico, tanto meno da un sistema assistito da un modello linguistico di grandi dimensioni (\gls{llm}): alcuni criteri richiedono un giudizio di merito sul contenuto (ad esempio la pertinenza di un testo alternativo rispetto al contesto informativo dell'immagine) che un \gls{llm} può supportare meglio di un controllo puramente sintattico, mentre altri richiedono l'osservazione runtime di comportamenti dinamici (ad esempio la gestione del focus della tastiera in un'interazione complessa) che restano più adatti a strumenti di test automatizzato tradizionali. La selezione dei criteri da includere nell'estensione è stata quindi condotta incrociando due fonti: da un lato il report WebAIM Million, che rileva annualmente su un campione di un milione di pagine web le violazioni \gls{wcag} più frequenti sul web reale, individuando così quali criteri concentrano il maggior numero di barriere effettive; dall'altro una valutazione di integrabilità tecnica con un'architettura basata su \gls{llm}, per identificare quali di quei criteri più problematici siano anche verificabili con un livello di affidabilità adeguato da un modello linguistico.

A questo doppio filtro --- prevalenza empirica delle violazioni e verificabilità tramite \gls{llm} --- si è aggiunto un terzo criterio, coerente con quanto argomentato nella sezione~\ref{subsec:cose-accessibilita}: a parità di frequenza e di fattibilità tecnica, sono stati preferiti i criteri capaci di intercettare la più ampia gamma di utenza, includendo non solo le disabilità permanenti ma anche le condizioni temporanee e situazionali. È la logica dei quattro pilastri POUR applicata in chiave selettiva: non si tratta di scegliere i criteri più facili da implementare, ma quelli che, a fronte di uno sforzo di verifica automatica limitato, producono il beneficio più ampio e più diffuso per chi naviga il web in condizioni di svantaggio, permanente o contingente che sia. I criteri specifici selezionati, e la metodologia con cui sono stati individuati a partire dai dati del WebAIM Million, sono oggetto della sezione~\ref{sec:selezione-criteri}.


\subsection{L'accessibilità come metodo di progettazione, non come rifinitura finale}
\label{subsec:accessibilita-by-design}

Quanto ricostruito nelle sezioni precedenti permette di formulare un'ultima considerazione, che ha un valore più metodologico che normativo. L'accessibilità, così come emerge dalle definizioni esaminate nella sezione~\ref{subsec:cose-accessibilita} e dalla logica dei quattro principi POUR discussa nella sezione~\ref{subsec:quattro-principi}, non è un insieme di regole da imparare a memoria né una rifinitura da applicare a un prodotto digitale già completato. È, piuttosto, un modo di progettare e sviluppare prodotti digitali migliori: più inclusivi, perché ampliano la platea di persone in grado di utilizzarli, e più efficaci, perché --- come si è visto a proposito del meccanismo \enquote{essenziale per una categoria di utenti, vantaggioso per molte altre} che ricorre nella norma EN~301~549 --- le soluzioni pensate per garantire l'accesso a chi si trova in una condizione di svantaggio finiscono quasi sempre per migliorare l'esperienza d'uso anche per chi quella condizione non ce l'ha. Le stesse linee guida \gls{agid}, del resto, distinguono esplicitamente un approccio \enquote{by design}, che integra l'accessibilità in ogni fase del ciclo di vita del prodotto fin dalla progettazione, da un approccio correttivo che interviene solo a posteriori, e indicano il primo come la strada più efficace, oltre che meno onerosa nel tempo\footcite[§~6.2]{agid:lg-82-2022}.

Tra il principio e la pratica, tuttavia, resta una distanza che i dati aiutano a misurare. Il WebAIM Million, un'analisi automatizzata condotta annualmente dal 2019 sulle home page di un milione di siti web tra i più visitati al mondo, fornisce la serie storica più ampia disponibile sullo stato reale dell'accessibilità del web. Nell'edizione più recente, relativa al 2026, il 95,9\% delle home page analizzate presentava almeno un errore rilevabile automaticamente e riconducibile con elevata probabilità a una violazione dei criteri di conformità \gls{wcag} 2.2 di livello A o AA, con una media di 56 errori per pagina\footcite{webaim:million2026}. Si tratta di un dato sostanzialmente stabile nel tempo: la quota di home page con almeno un errore rilevato era già al 97,8\% nella prima edizione del 2019 e, pur con oscillazioni minime anno su anno, non è mai scesa sotto il 94,8\%, toccato nel 2025. In altre parole, per l'intero periodo in cui questo tipo di monitoraggio è stato condotto, la stragrande maggioranza dei siti web più visitati al mondo ha continuato a presentare barriere di accessibilità rilevabili in modo automatico --- un segnale che l'approccio \enquote{by design} resta, nella pratica corrente dello sviluppo web, l'eccezione più che la norma.

Questo scarto tra principio e pratica ha un'implicazione diretta per come affrontare il problema. Se la maggior parte del web esistente non è stata progettata secondo criteri di accessibilità, non è sufficiente promuovere l'adozione di questi criteri solo per i contenuti che verranno creati in futuro: è necessario intervenire su due fronti contemporaneamente. Da un lato, incoraggiare che i nuovi contenuti siano progettati fin dall'origine in ottica di accessibilità, così da non riprodurre gli stessi errori sistematici che il WebAIM Million rileva costantemente. Dall'altro, rendere possibile ed economicamente sostenibile la verifica e il miglioramento di ciò che è già stato pubblicato, affinché il numero più ampio possibile di utenti --- comprese le persone con disabilità permanenti, temporanee e situazionali di cui si è discusso nella sezione~\ref{subsec:cose-accessibilita} --- possa beneficiarne. È in questo secondo fronte, quello della verifica di contenuti già esistenti su larga scala, che si colloca il problema pratico affrontato nella sezione successiva.


\section{Il problema pratico: la persistenza delle barriere di accessibilità}
\label{sec:problema-pratico}

% TODO: approfondimento dei dati WebAIM Million (serie storica 2019-2026) sulle
% tipologie di errore più frequenti riscontrate sul web reale, come base empirica
% per la selezione dei criteri WCAG.


\section{Selezione dei criteri WCAG oggetto di verifica}
\label{sec:selezione-criteri}

% TODO: illustrazione dei tre criteri incrociati per selezionare il sottoinsieme di
% criteri WCAG 2.1 AA effettivamente implementati nell'estensione (prevalenza
% empirica, integrabilità con LLM, copertura della più ampia gamma di utenza),
% con elenco motivato dei criteri selezionati.


\section{Analisi dei requisiti}
\label{sec:analisi-requisiti}

% TODO: requisiti funzionali e non funzionali dell'estensione, casi d'uso
% principali, eventuale analisi preventiva dei rischi.


\section{Progettazione}
\label{sec:progettazione}

% TODO: architettura dell'estensione browser, scelte tecnologiche, modalità di
% integrazione con il modello linguistico, design pattern adottati.


\section{Implementazione}
\label{sec:implementazione}

% TODO: aspetti realizzativi più significativi, con frammenti di codice
% commentati dove utile.


\section{Verifica, validazione e risultati}
\label{sec:verifica-risultati}

% TODO: strategia di test ed esito; risultati qualitativi (prodotto visto dal
% lato dell'utente finale) e quantitativi (copertura requisiti, copertura test,
% metriche di codice, accuratezza della verifica automatica).
